\section{Introduction}
\label{sec:intro}

Predicting the toxicological profiles of novel chemical entities in silico is a foundational challenge in computational pharmacology, environmental hazard screening, and modern regulatory safety~\cite{huang2016tox21}. In vivo toxicology assays are prohibitively slow, resource-intensive, and ethically constrained, creating an urgent imperative for highly accurate, generalizable machine learning algorithms capable of screening vast combinatorial chemical libraries. The Toxicology in the 21st Century (Tox21) data challenge formalizes this task as a 12-task multi-label classification problem across major nuclear receptor (NR) and stress response (SR) pathways~\cite{huang2016tox21}.

Historically, molecular machine learning has relied predominantly on fixed, discrete molecular representations, most notably circular Extended-Connectivity Fingerprints (ECFP/Morgan fingerprints) and handcrafted physicochemical descriptor sets~\cite{wu2018moleculenet}. While computationally efficient and empirically robust when paired with tree ensembles such as Random Forest and XGBoost, fixed circular fingerprints suffer from fundamental theoretical limitations. By hashing local atom environments into fixed-width bit vectors, fingerprints induce bit collisions, ignore continuous spatial and geometric variations, and treat all chemical subgraphs as rigid, unweighted components regardless of biological context.

Graph Neural Networks (GNNs) offer an expressive paradigm that resolves these representational constraints by treating molecules as natural mathematical graphs, where atoms correspond to attributed nodes and chemical bonds correspond to directed edges~\cite{duvenaud2015convolutional,gilmer2017neural}. Through iterative neural message passing, node representations continuously aggregate contextual information from topological neighborhoods, enabling the model to learn flexible, task-specific molecular representations end-to-end.

In this Phase~2 investigation, we conduct a disciplined, leak-free empirical evaluation comparing graph deep learning architectures against frozen classical machine learning baselines established in Phase~1. We specifically formulate:
\begin{enumerate}
    \item A Message Passing Neural Network (MPNN) featuring continuous edge networks, gated recurrent unit (GRU) node state updates, and an order-invariant Set2Set attention pooling readout~\cite{gilmer2017neural,vinyals2015order}.
    \item A multi-modal hybrid architecture (\textbf{H1}: \texttt{mpnn+desc}) that fuses learned graph representations with 213 standardized 2D RDKit physicochemical descriptors before classification.
    \item Two strict experimental controls: an \textbf{FP-MLP control} (\texttt{fpmlp+both}) to decouple neural architecture capacity from graph topology learning, and a \textbf{shuffled-descriptor control} to prove that hybrid gains stem from genuine molecular physics rather than parameter regularization.
    \item A late-fusion rank-average ensemble (\textbf{H2}: \texttt{ens\_xgb\_mpnn+desc}) combining gradient-boosted decision trees and message passing networks to exploit complementary inductive biases.
\end{enumerate}

All evaluations are conducted under a strict out-of-distribution Bemis--Murcko scaffold split~\cite{bemis1996properties} with three random seeds, masked multi-task cross-entropy loss, and hyperparameter configurations frozen prior to test evaluation. Our results demonstrate that while tree ensembles on physicochemical descriptors remain exceptionally competitive on small benchmarks, graph message passing provides superior structural representations over flat MLPs, and late-fusion ensembling delivers consistent, statistically significant improvements across all 12 toxicity endpoints ($p = 0.00049$).
